% !TeX root=../main.tex
% This section is included from tex/chapter4.tex.
\section{آزمایش مشاهده‌ای هدایت گشت به مراکز هدف}
\label{sec:ch4-center-target-observation}

\subsection{انگیزه و تعریف مرکز هدف}
\label{subsec:ch4-center-observation-motivation}
در فصل سوم مراکز کلاس‌ها به‌صورت میانگین توزیع‌های همسایگی تعریف شدند. بهبود جدایی این مراکز، بدون توجه به پراکندگی نمونه‌ها در اطراف آن‌ها، لزوماً به بهبود تصمیم طبقه‌بندی منجر نمی‌شود. برای مشاهده تجربی این موضوع، ابتدا توزیع‌هایی را به‌عنوان \emph{مراکز هدف} تعیین می‌کنیم و سپس با انتخاب احتمالی ویژگی‌ها تلاش می‌کنیم بازنمایی نمونه‌ها را به آن توزیع‌ها نزدیک کنیم. مراکز هدف را با $U_c^{(t)}$ و مراکز \emph{تحقق‌یافته} از گشت‌های آموزشی را با $\widehat{\mathcal V}_c^{(t)}$ نشان می‌دهیم. این دو مفهوم یکسان نیستند: حتی بهترین سیاست محلی ممکن است نتواند مرکز هدف را ایجاد کند.

در این آزمایش سیاست‌ها تابع کلاس هدف و وضعیت جاری گشت هستند؛ هدف آن \emph{مشاهده رفتار معیارها} است، نه آموزش یک ماتریس سیاست مشترک یا اثبات همگرایی یک الگوریتم یادگیری.

\subsection{سیاست انتخاب ویژگی با استفاده از توزیع یک‌گامی}
\label{subsec:ch4-center-observation-local-policy}
فرض کنیم گشت از نمونه $X_s$ آغاز شده، در نمونه $X_j$ قرار دارد و هدف آن رسیدن به توزیع کلاس $c$ است. برای ویژگی مجاز $b$، توزیع دقیق برچسب مقصد را روی گراف محدود آموزش با $Q_{j,b}^{(s)}$ نمایش می‌دهیم. مجموعه مقاصد در محاسبه این توزیع، نمونه جاری و نمونه آغازین را شامل نمی‌شود. اگر $B_j^{(s)}$ مجموعه ویژگی‌های مجاز در وضعیت جاری باشد، سیاست محلی از رابطه زیر به دست می‌آید:
\begin{equation}\label{eq:ch4-center-observation-policy}
\boxed{
\pi_{j\mid s,c}^{(t)}\in
\underset{\pi\in\Delta(B_j^{(s)})}{\arg\min}\;
\operatorname{JSD}\left(
\sum_{b\in B_j^{(s)}}\pi_bQ_{j,b}^{(s)},\;U_c^{(t)}\right).
}
\end{equation}
در نمایش ماتریسی، عبارت نخست داخل واگرایی برابر $(Q_j^{(s)})^\top\pi$ است. به دلیل محدب‌بودن JSD نسبت به توزیع نخست و خطی‌بودن ترکیب توزیع‌های یک‌گامی، مسئله نسبت به $\pi$ محدب است. برای حل عددی آن از روش گرادیان شرطی (\lr{Frank--Wolfe}) با جست‌وجوی خطی تقریبی استفاده می‌کنیم. پاسخ محاسبه‌شده، بهینه‌سازی \emph{یک‌گامی} نسبت به مرکز هدف است و بهینگی گشت کامل را تضمین نمی‌کند.

این سیاست نه فقط در مبدأ، بلکه در \emph{تمام گره‌های واسط} برای وضعیت جاری و با حذف بازگشت به نمونه آغازین محاسبه می‌شود. سیاست آموزشی از کلاس معلوم نمونه مبدأ به‌عنوان کلاس هدف استفاده می‌کند. در اعتبارسنجی، چون برچسب نمونه ناشناخته است، به ازای هر کلاس فرضی گشت مستقلی اجرا می‌شود و برچسب واقعی فقط پس از تصمیم نهایی برای ارزیابی به کار می‌رود.

\subsection{تولید هدف‌های متوالی و محاسبه مراکز واقعی}
\label{subsec:ch4-center-observation-iteration}
در شروع، برای ساخت هدف‌هایی با جدایی زیاد از توزیع‌های نزدیک به بردارهای یک‌داغ استفاده می‌کنیم:
\begin{equation}\label{eq:ch4-center-observation-initial-target}
U_c^{(0)}=(1-\varepsilon)e_c+\varepsilon\boldsymbol u,\qquad
\boldsymbol u=(1/K,\ldots,1/K),\qquad\varepsilon=0.08.
\end{equation}
پس از محاسبه سیاست‌های محلی بر اساس $U^{(t)}$، گشت‌های تصادفی چندگامی اجرا می‌شوند. هر گشت با ضرایب $\lambda$ و مجموع وزن همان گشت نرمال می‌شود و میانگین گشت‌ها، بازنمایی $\widehat{\mathscr Y}_s^{(t)}$ را می‌سازد. مراکز واقعی با همان منطق میانگین‌گیری فصل سوم محاسبه می‌شوند:
\begin{equation}\label{eq:ch4-center-observation-realized}
\boxed{
\widehat{\mathcal V}_c^{(t)}=
\frac{1}{|C_c^{\mathrm{train}}|}
\sum_{s\in C_c^{\mathrm{train}}}\widehat{\mathscr Y}_s^{(t)}.
}
\end{equation}
پس از \emph{ارزیابی با مراکز واقعی همین تکرار}، برای تکرار بعد هدف‌ها را با رابطه زیر تعیین می‌کنیم:
\begin{equation}\label{eq:ch4-center-observation-update}
\boxed{U_c^{(t+1)}=(1-\eta)U_c^{(t)}+
\eta\widehat{\mathcal V}_c^{(t)},\qquad\eta=0.5.}
\end{equation}
این رابطه یک تکرار مشاهده‌ای برای سازگارکردن مرکز هدف با خروجی گشت است؛ هیچ جمله‌ای برای افزایش مستقل جدایی کلاس‌ها یا کمینه‌سازی خطای تصمیم ندارد. کوچک‌شدن فاصله $U^{(t)}$ و $\widehat{\mathcal V}^{(t)}$ به‌تنهایی به معنای بهبود طبقه‌بندی یا تضمین همگرایی سیاست نیست.

تنظیمات آزمایش شامل سه کلاس، پنج ویژگی، $300$ نمونه آموزشی و $90$ نمونه اعتبارسنجی است. در هر مبدأ $100$ گشت با $\beta=0.4$، $\lambda=0.8$ و سقف $12$ گام اجرا شده است. بهینه‌ساز محلی حداکثر $45$ تکرار و بذر گشت $20261104$ دارد. $49$ وضعیت از تکرار صفر تا $48$ ثبت شده است؛ دقت اعتبارسنجی در $25$ تکرار واقعاً اندازه‌گیری شده و برای سایر تکرارها برآورد نشده است. بذرهای گشت در تکرارها ثابت نگه داشته شدند تا بخش مهمی از نوسان مونت‌کارلو حذف شود.

\subsection{پروتکل تصمیم و سه وضعیت منتخب}
\label{subsec:ch4-center-observation-cases}
در سیاست یکنواخت، برای هر نمونه اعتبارسنجی فقط یک بازنمایی مستقل از کلاس فرضی ساخته و با مراکز آموزش همان سیاست مقایسه می‌شود. در سیاست هدایت‌شده، برای هر فرض کلاس $c$ بازنمایی جداگانه $\widehat{\mathscr Y}_i^{(t,c)}$ ایجاد می‌شود و قاعده تصمیم زیر اعمال می‌گردد:
\begin{equation}\label{eq:ch4-center-observation-decision}
\widehat y_i^{(t)}=
\underset{c}{\arg\min}\;
\operatorname{JSD}\left(
\widehat{\mathscr Y}_i^{(t,c)},\widehat{\mathcal V}_c^{(t)}\right).
\end{equation}
در هر دو روش، مراکز از داده آموزشی به دست می‌آیند و برچسب‌های اعتبارسنجی در ساخت گشت وارد نمی‌شوند. بااین‌حال، \emph{پروتکل‌های استنتاج یکنواخت و هدایت‌شده یکسان نیستند} و تفاوت دقت این دو را نباید صرفاً به تغییر سیاست نسبت داد.

سه وضعیت برای مقایسه انتخاب می‌شوند: سیاست یکنواخت؛ نخستین اجرای هدایت‌شده با هدف تقریباً یک‌داغ و بدون اصلاح قبلی آن؛ و اجرای پنجم، یعنی تکرار چهار، که در میان $25$ ارزیابی انجام‌شده بیشترین دقت متوازن را ثبت کرده است. هدف و مرکز تحقق‌یافته در دو وضعیت هدایت‌شده به‌ترتیب چنین‌اند:
\begin{equation}\label{eq:ch4-center-observation-target-first}
U^{(0)}=\begin{pmatrix}
.9467&.0267&.0267\\ .0267&.9467&.0267\\ .0267&.0267&.9467
\end{pmatrix},\qquad
\widehat{\mathcal V}^{(0)}=\begin{pmatrix}
.5874&.3551&.0575\\ .1497&.7122&.1381\\ .0383&.2905&.6712
\end{pmatrix}.
\end{equation}
\begin{equation}\label{eq:ch4-center-observation-target-best}
U^{(4)}=\begin{pmatrix}
.5393&.3929&.0678\\ .1577&.6915&.1508\\ .0556&.3206&.6237
\end{pmatrix},\qquad
\widehat{\mathcal V}^{(4)}=\begin{pmatrix}
.4629&.4540&.0831\\ .1857&.6396&.1747\\ .0762&.3713&.5525
\end{pmatrix}.
\end{equation}
مراکز واقعی سیاست یکنواخت عبارت‌اند از:
\begin{equation}\label{eq:ch4-center-observation-uniform}
\widehat{\mathcal V}^{\mathrm{uni}}=\begin{pmatrix}
.3249&.4701&.2050\\ .2215&.5412&.2373\\ .1812&.4475&.3714
\end{pmatrix}.
\end{equation}
بهترین بودن اجرای پنجم، تنها به معنای بیشترین دقت متوازن \emph{اندازه‌گیری‌شده در همین مسیر و همین مجموعه اعتبارسنجی} است؛ انتخاب پس‌نگر روی اعتبارسنجی جایگزین ارزیابی نهایی با مجموعه آزمون نیست.

\subsection{معیارهای تراکم، درست‌نمایی و جدایی}
\label{subsec:ch4-center-observation-metrics}
برای هر مرکز واقعی، معیارهای $T$، $BT$، $TNL$، $BTNL$ و $I$ را طبق فصل سوم محاسبه می‌کنیم. همچنین برای زوج‌مرکزها میانگین JSD، میانگین و کمینه $\sqrt{\operatorname{JSD}}$ و میانگین شباهت نرمال‌شده JSD را گزارش می‌دهیم. به دلیل رابطه $s(P,Q)=1-\operatorname{JSD}(P,Q)/\log 2$، میانگین شباهت و میانگین JSD دو نمایش خطی از یک اطلاعات هستند، ولی میانگین جذر JSD را نباید با جذر میانگین JSD اشتباه گرفت.

برای سنجش فشردگی درون‌کلاسی، افزون بر میانگین JSD، دو معیار فاصله و شباهت زیر را برای بازنمایی‌های آموزش تعریف می‌کنیم:
\begin{equation}\label{eq:ch4-center-observation-JS-within}
\begin{aligned}
R_{\mathrm{JS}}&=\frac1K\sum_c\frac1{n_c}\sum_{i\in C_c}
\operatorname{JSD}(\mathscr Y_i,\mathcal V_c),\\
R_{\sqrt{\mathrm{JS}}}&=\frac1K\sum_c\frac1{n_c}\sum_{i\in C_c}
\sqrt{\operatorname{JSD}(\mathscr Y_i,\mathcal V_c)},\\
S_{\mathrm{JS}}&=\frac1K\sum_c\frac1{n_c}\sum_{i\in C_c}
s(\mathscr Y_i,\mathcal V_c).
\end{aligned}
\end{equation}
این کمیت‌ها مکمل شعاع اصلی $R$ فصل سوم، یعنی میانگین $D_{\mathrm{KL}}(\mathscr Y_i\|\mathcal V_c)$ هستند. چون مرکز، میانگین حسابی بازنمایی‌های کلاس است، برای هر کلاس داریم:
\begin{equation}\label{eq:ch4-center-observation-R-likelihood}
\begin{aligned}
R_c&=H(\mathcal V_c)-\frac1{n_c}\sum_{i\in C_c}H(\mathscr Y_i),\\
\frac1{n_c}\sum_{i\in C_c}\sum_j\mathscr Y_i(j)\log\mathcal V_c(j)
&=-H(\mathcal V_c).
\end{aligned}
\end{equation}
رابطه دوم، \emph{میانگین لگاریتم درست‌نمایی مشاهدات همسایگی تحت مدل مرکز} است، نه احتمال درست‌بودن کلاس نمونه. در نتیجه، درست‌نمایی مرکز، پراکندگی $R$ و پراکندگی مبتنی بر JSD باید جدا از هم تفسیر شوند.

\begin{table}[htbp]
\centering\small
\caption{معیارهای مبتنی بر مراکز \emph{واقعی} برای سه وضعیت منتخب؛ لگاریتم‌ها طبیعی‌اند.}
\label{tab:ch4-center-observation-geometry}
\begin{tabular}{rccc}\hline
معیار & یکنواخت & هدایت‌شده، صفر & هدایت‌شده، چهار\\\hline
$I$ & .01818 & .26873 & .14111\\
$T=BT$ & .41247 & .65694 & .55166\\
$TNL=BTNL$ & .90963 & .42337 & .60350\\
میانگین JSD زوج‌مرکز & .01336 & .18860 & .10143\\
کمینه جذر JSD زوج‌مرکز & .08223 & .32937 & .21640\\
میانگین جذر JSD زوج‌مرکز & .11202 & .42477 & .30709\\
میانگین شباهت JSD مراکز & .98072 & .72790 & .85367\\\hline
\end{tabular}
\end{table}
برابری‌های $T=BT$ و $TNL=BTNL$ در این جدول، تنها به این دلیل رخ داده‌اند که انتساب هم‌نام کلاس‌ها در این سه ماتریس بهترین انتساب جایگشتی است؛ این برابری عمومی نیست.

\begin{table}[htbp]
\centering\small
\caption{تراکم بازنمایی‌های \emph{آموزشی}، میانگین لگاریتم درست‌نمایی مرکز و معیارهای زیان. مقادیر کوچک‌تر برای فاصله‌ها و NLL مطلوب‌اند.}
\label{tab:ch4-center-observation-within}
\begin{tabular}{rccc}\hline
معیار & یکنواخت & هدایت‌شده، صفر & هدایت‌شده، چهار\\\hline
$R$ مبتنی بر KL & .01542 & .05466 & .01513\\
$R_{\mathrm{JS}}$ & .00388 & .01326 & .00376\\
$R_{\sqrt{\mathrm{JS}}}$ & .05508 & .09995 & .05295\\
$S_{\mathrm{JS}}$ & .99440 & .98087 & .99457\\
لگاریتم درست‌نمایی مرکز & -1.02989 & -.79850 & -.90567\\
NLL ساده آموزش & 1.06341 & .65381 & .84258\\
NLL نرمال‌شده آموزش & .65689 & .39452 & .40615\\
معیار حاشیه‌ای $J$ آموزش & .87020 & .28292 & .18257\\\hline
\end{tabular}
\end{table}
دقت آموزش در دو ستون هدایت‌شده با استفاده از \emph{برچسب واقعی مبدأ در ساخت سیاست} به دست آمده و یک معیار \emph{شرطی و اوراکلی} است؛ بنابراین آن را معادل دقت آموزش طبقه‌بند معمولی تلقی نمی‌کنیم.

\begin{table}[htbp]
\centering\small
\caption{ارزیابی اعتبارسنجی با مراکز واقعی هر سیاست؛ روش‌های یکنواخت و هدایت‌شده پروتکل‌های استنتاج متفاوت دارند.}
\label{tab:ch4-center-observation-validation}
\begin{tabular}{rccc}\hline
معیار & یکنواخت & هدایت‌شده، صفر & هدایت‌شده، چهار\\\hline
دقت متوازن (درصد) & 84.86 & 72.61 & 83.83\\
دقت معمول (درصد) & 81.11 & 70.00 & 83.33\\
KL تا مرکز فرضیه درست & .01280 & .04664 & .01126\\
JSD تا مرکز فرضیه درست & .00323 & .01108 & .00286\\
میانگین جذر JSD تا مرکز درست & .05110 & .09086 & .04620\\
میانگین شباهت JSD تا مرکز درست & .99533 & .98401 & .99587\\
NLL ساده مبتنی بر KL & 1.05872 & 1.00289 & 1.04910\\
NLL نرمال‌شده مبتنی بر KL & .47910 & .70376 & .60724\\
NLL مبتنی بر شباهت JSD & 1.08425 & 1.06968 & 1.08217\\\hline
\end{tabular}
\end{table}
فاصله نمونه اعتبارسنجی تا «مرکز فرضیه درست» صرفاً در مرحله \emph{محاسبه معیار} از برچسب واقعی استفاده می‌کند؛ تولید بازنمایی‌های مربوط به فرضیه‌های مختلف از این برچسب بی‌اطلاع است. همچنین مقدار NLL در این بخش بر اساس امتیازهای نامساوی هر فرضیه و با وزن یکسان کلاس‌ها محاسبه شده است؛ به دلیل وابسته‌بودن گشت هدایت‌شده به کلاس فرضی، نباید آن را همان NLL یک بازنمایی مشترک در فصل سوم فرض کرد.

\subsection{تفسیر مسیر و محدودیت شاخص‌های منفرد}
\label{subsec:ch4-center-observation-interpretation}
از اجرای اول تا اجرای پنجم هدایت‌شده، $I$ از $.26873$ به $.14111$ کاهش یافته است، اما $R$ از $.05466$ به $.01513$ و میانگین جذر JSD درون‌کلاسی از $.09995$ به $.05295$ رسیده است. هم‌زمان دقت متوازن اعتبارسنجی از $72.61$ به $83.83$ درصد افزایش یافته است. بنابراین افزایش جدایی مراکز \emph{به‌تنهایی} توضیح‌دهنده دقت نیست و تغییر پراکندگی نمونه‌ها باید همراه آن سنجیده شود. به‌علاوه، NLL ساده اعتبارسنجی میان این دو اجرای هدایت‌شده از $1.00289$ به $1.04910$ افزایش یافته، درحالی‌که دقت متوازن بیشتر شده است. این مشاهده یک شاهد محدود و مشخص از هم‌راستا نبودن کامل NLL ساده با دقت است، نه رد عمومی این زیان.

در تکرارهای دیرتر، کاهش اختلاف مراکز هدف و واقعی ادامه یافته ولی جدایی مراکز تقریباً از بین رفته است: در تکرار $48$، $I\simeq0.000794$ و دقت متوازن اندازه‌گیری‌شده حدود $34.04$ درصد، نزدیک به مبنای $33.33$ درصد در مسئله سه‌کلاسه، بوده است. این روند نه اثبات همگرایی به یک نقطه یکتا، بلکه مشاهده‌ای از امکان فروپاشی تفکیک در روش سازگارسازی هدف با مرکز است. از $49$ تکرار تنها $25$ تکرار ارزیابی اعتبارسنجی دارند؛ مقادیر غایب در تحلیل درون‌یابی نشده‌اند.

در مجموع، هیچ‌یک از معیارهای صرفاً وابسته به مراکز، تراکم مبتنی بر KL یا JSD، و NLL ساده به‌تنهایی تضمین‌کننده بهبود دقت طبقه‌بندی نیستند. این نتیجه، انگیزه‌ای تجربی برای طراحی تابع هدفی است که \emph{جدایی مراکز و پراکندگی/حاشیه تصمیم نمونه‌ها را هم‌زمان} در نظر بگیرد. تعریف نظری دقیق معیارهای مکمل مبتنی بر فاصله و شباهت JSD می‌تواند در تکمیل فصل سوم صورت گیرد؛ در اینجا تنها محاسبه و نقش تجربی آن‌ها بررسی شده است.

کد محاسبه سیاست، داده‌های بازنمایی نمونه‌ها، معیارهای تمامی مراکز $49$ تکرار و خروجی‌های تفصیلی سه وضعیت در مخزن پژوهش در مسیر \lr{\texttt{code/experiments/center\_matching\_comparison}} ثبت شده‌اند. این نتایج متعلق به یک سناریوی مصنوعی و یک بذرند؛ نقطه بیشینه دقت با اعتبارسنجی انتخاب شده و مجموعه آزمون برای انتخاب آن استفاده نشده است.
